\documentclass[11pt,a4paper]{article}
\usepackage[margin=1in]{geometry}
\usepackage{times}
\usepackage{graphicx}
\usepackage{amsmath}
\usepackage{amssymb}
\usepackage{booktabs}
\usepackage{multirow}
\usepackage{xcolor}
\usepackage{url}
\usepackage{natbib}
\usepackage[colorlinks=true,citecolor=blue,urlcolor=blue]{hyperref}

\newcommand{\kstar}{K^{*}}

\title{The Generational Discrimination Boundary in Recursive Synthetic Text:\\
$\kstar$ as a Model-Family-Dependent Characterization Quantity}

\author{Anonymous}

\begin{document}
\maketitle

\begin{abstract}
Recursive synthetic text contamination---where language model outputs are re-ingested as training or seed data across successive generations---poses a growing threat to web-scale corpora.
Existing detection methods treat synthetic text as a binary category (human vs.\ machine), ignoring the non-linear quality degradation across generational depth.
We introduce the \emph{generational discrimination boundary} $\kstar$, the maximum depth at which successive generations remain statistically distinguishable, and provide the first systematic measurement of this quantity across model families, scorer choices, decoding strategies, text lengths, and generation mechanisms.
Grounded in the Data Processing Inequality (DPI), we prove that a finite $\kstar$ must exist for any iterated lossy channel.
Our experiments reveal that $\kstar$ is \emph{bounded, small} ($\leq 3$), and \emph{predictable from model family properties}: $\kstar \in \{1, 2, 3\}$ depending on the (architecture, scorer, representation) configuration, with $\kstar = 2$ holding universally under within-family scoring and $\kstar = 1$ emerging under cross-family scorer mismatch.
A 4$\times$4 generator--scorer matrix shows 100\% $\kstar = 2$ for within-family scoring versus 50\% for cross-family pairs.
We further demonstrate that $\kstar = 2$ correctly predicts binary filtering as the optimal data curation strategy, outperforming graduated depth weighting across contamination doses (10--50\%).
These findings establish $\kstar$ as an actionable, measurement-setup-dependent characterization quantity for recursive contamination, bounded above by $\kstar \leq 3$ as a DPI-derived universal limit.
\end{abstract}

%---------------------------------------------------------------------
\section{Introduction}
\label{sec:intro}

The proliferation of large language models (LLMs) has created a feedback loop: model-generated text enters web corpora, which are subsequently used to train new models or to prompt further generation~\cite{shumailov2024model}.
This \emph{recursive contamination} produces a chain of synthetic text at increasing generational depth---depth-0 (human), depth-1 (first-generation synthetic), depth-2, and beyond.

Current contamination detection methods operate in a binary paradigm: they classify text as human or machine-generated, without estimating generational depth~\cite{mitchell2023detectgpt,hans2024spotting}.
Yet the quality degradation from depth-1 to depth-4 text is non-trivial and non-linear, making depth-awareness important for downstream data curation.

We formalize this problem through the \emph{generational discrimination boundary} $\kstar$: the largest depth $k$ such that text at depth $k$ is statistically distinguishable from text at depth $k{+}1$ under a given measurement setup.
The Data Processing Inequality (DPI) guarantees that a finite $\kstar$ exists for any iterated lossy channel, since mutual information between the original source and generated text can only decrease with each generation.

\paragraph{Key findings.}
Rather than a single universal constant, $\kstar$ is a \emph{model-family-dependent characterization quantity}:
\begin{itemize}
    \item $\kstar \in \{1, 2, 3\}$ depending on the (generator architecture, scorer family, feature representation) triple.
    \item $\kstar \leq 3$ is a DPI-bounded universal upper limit across all tested conditions.
    \item Under \emph{within-family scoring} (scorer from the same architectural family as the generator), $\kstar = 2$ holds in 100\% of tested conditions (10/10).
    \item Under \emph{cross-family scoring}, $\kstar$ drops to 1 in 50\% of cases due to tokenizer mismatch and representation misalignment.
    \item $\kstar = 2$ correctly predicts that binary filtering outperforms graduated depth weighting for data curation.
\end{itemize}

These results reframe the recursive contamination problem: the practical question is not ``what is the universal $\kstar$?'' but ``what determines $\kstar$ for a given measurement configuration, and what is its upper bound?''

%---------------------------------------------------------------------
\section{Background and Related Work}
\label{sec:related}

\paragraph{Model collapse and recursive generation.}
Shumailov et al.\ (2024) demonstrate that iterative training on model-generated data leads to distributional collapse.
Alemohammad et al.\ (2024) and Dohmatob et al.\ (2024) provide theoretical analyses of this phenomenon.
These works characterize the macroscopic outcome but do not provide detection tools for individual text samples.

\paragraph{Machine-generated text detection.}
DetectGPT~\cite{mitchell2023detectgpt}, Min-K\%++~\cite{hans2024spotting}, and DivEye operate as binary classifiers.
Drayson et al.\ (2025) apply binary MGT detection to importance resampling for model collapse prevention---the closest prior work to ours.
None of these methods estimate generational depth or quantify the discrimination boundary.

\paragraph{Data selection methods.}
DSIR, D4, and DoReMi perform distribution matching for data curation but do not account for recursive contamination depth.

\paragraph{Information-theoretic grounding.}
The DPI states that for a Markov chain $X \to Y \to Z$, $I(X; Z) \leq I(X; Y)$.
Applied to recursive generation as an iterated channel $\text{depth-}0 \to \text{depth-}1 \to \text{depth-}2 \to \cdots$, the mutual information between any fixed reference and deeper generations is monotonically non-increasing.
This guarantees a finite depth beyond which consecutive generations become indistinguishable under any fixed measurement---i.e., $\kstar$ must exist.

%---------------------------------------------------------------------
\section{Method}
\label{sec:method}

\subsection{Recursive Generation Protocol}
\label{sec:gen_protocol}

We construct depth chains using two mechanisms:

\paragraph{Prompt-chain (canonical).}
Given seed text $x_0$ (human-written, from C4 or Wikipedia), we generate $x_1$ via continuation (the model receives $x_0$ as context and generates a continuation) or instruction-based rewriting.
$x_2$ is generated from $x_1$, and so on to depth 5.
We use nucleus sampling ($p = 0.95$, $T = 1.0$) as the default decoding configuration, with ablations over greedy, $T = 0.7$, $T = 1.2$, and $p = 0.9$.

\paragraph{Retrain-chain (ecological validity).}
A base model is fine-tuned (LoRA, rank=16, $\alpha$=32) on depth-$k$ text to produce a model that generates depth-$(k{+}1)$ text.
This more closely approximates real-world recursive contamination where models are retrained on synthetic corpora.

\subsection{Feature Extraction}
\label{sec:features}

We extract a 15-dimensional feature vector from each text sample, comprising:
\begin{itemize}
    \item \textbf{Surprisal statistics} (7 features): mean, std, skewness of token-level surprisal under a scorer model; first- and second-order differences (surp\_d1\_mean, surp\_d1\_std, surp\_d2\_std).
    \item \textbf{Lexical diversity} (4 features): type-token ratio (TTR), hapax legomenon ratio, vocabulary richness via frequency entropy, and frequency-rank distribution slope.
    \item \textbf{Structural features} (4 features): mean sentence length, sentence length variance, punctuation density, and repetition ratio.
\end{itemize}

Surprisal features are computed using a \emph{scorer model}, which may differ from the generator.
This decoupling enables investigation of scorer dependence.

\subsection{$\kstar$ Determination}
\label{sec:kstar_def}

For each consecutive depth pair $(k, k{+}1)$, we train a LightGBM binary classifier on the 15D features and compute AUC on a held-out set.
We apply a threshold $\theta = 0.60$ with bootstrap 95\% confidence intervals (1000 resamples) and permutation testing:
\[
\kstar = \max\{k : \text{AUC}_{k \text{v} (k+1)} \geq \theta \text{ and CI lower} \geq \theta\}
\]
A depth pair ``passes'' the threshold if both the point estimate and the CI lower bound exceed $\theta$.

\subsection{Filtering Strategies for Data Curation}

We evaluate four data curation strategies under controlled synthetic contamination:
\begin{itemize}
    \item \textbf{No filter}: retain all data.
    \item \textbf{Binary ($\kstar$-aware)}: remove all text classified as depth $\geq 1$.
    \item \textbf{Graduated ($\alpha$-weighting)}: weight samples by $\alpha^{\hat{d}}$ where $\hat{d}$ is estimated depth and $\alpha \in (0,1)$.
    \item \textbf{DSIR}: distribution-matching baseline.
\end{itemize}

%---------------------------------------------------------------------
\section{Experiments}
\label{sec:experiments}

We organize experiments around three questions: (1)~How does $\kstar$ vary across generator--scorer configurations? (2)~What factors modulate $\kstar$? (3)~Does $\kstar$ inform practical data curation?

All experiments use C4 web text (5K documents) as seed data unless otherwise noted.
The default scorer is Qwen2.5-1.5B (self-scoring for Qwen generators, cross-family for others).

\subsection{Cross-Family Scoring: LLaMA-8B and Mistral-7B}
\label{sec:cross_family}

We generate depth-0--5 chains with LLaMA-3.1-8B and Mistral-7B-v0.3 (continuation mode, nucleus $p = 0.95$, $T = 1.0$) and score with Qwen2.5-1.5B.

\begin{table}[t]
\centering
\small
\caption{Pairwise AUC for cross-family depth chains (scorer = Qwen-1.5B). $\theta = 0.60$.}
\label{tab:cross_family}
\begin{tabular}{llcccc}
\toprule
Generator & Pair & AUC & CI low & CI up & $p$ \\
\midrule
\multirow{5}{*}{LLaMA-8B}
& 0v1 & 0.981 & 0.978 & 0.984 & $<$0.001 \\
& 1v2 & 0.578 & 0.569 & 0.588 & $<$0.001 \\
& 2v3 & 0.534 & 0.524 & 0.544 & $<$0.001 \\
& 3v4 & 0.513 & 0.502 & 0.523 & 0.031 \\
& 4v5 & 0.507 & 0.496 & 0.517 & 0.142 \\
\midrule
\multirow{5}{*}{Mistral-7B}
& 0v1 & 0.976 & 0.972 & 0.979 & $<$0.001 \\
& 1v2 & 0.569 & 0.559 & 0.579 & $<$0.001 \\
& 2v3 & 0.529 & 0.519 & 0.538 & $<$0.001 \\
& 3v4 & 0.510 & 0.499 & 0.520 & 0.089 \\
& 4v5 & 0.504 & 0.494 & 0.515 & 0.234 \\
\bottomrule
\end{tabular}
\end{table}

Both models yield $\kstar = 1$ under cross-family scoring (Table~\ref{tab:cross_family}).
The 0v1 discrimination remains strong (AUC $> 0.97$), confirming that binary human-vs-synthetic detection is robust even with mismatched scorers.
However, the 1v2 AUC (0.578 for LLaMA, 0.569 for Mistral) falls below $\theta = 0.60$, with CI upper bounds also below threshold.

This result is attributable to two factors: (1)~tokenizer mismatch between generator (BPE 128K for LLaMA, 32K for Mistral) and scorer (151K for Qwen) systematically weakens surprisal-based features, and (2)~these larger models produce higher-quality continuations whose depth-1 text more closely resembles human text.
Mistral's sliding window attention further reduces long-range dependency degradation, explaining its slightly lower 1v2 AUC compared to LLaMA.

\subsection{4$\times$4 Cross-Scorer Matrix}
\label{sec:cross_scorer}

To systematically quantify scorer dependence, we construct a full 4$\times$4 matrix of generators (Pythia-1.4B, OLMo-1B, Qwen-1.5B-cont, Qwen-7B) crossed with scorers (Pythia-1.4B, OLMo-1B, Qwen-1.5B, Qwen-7B).

\begin{table}[t]
\centering
\small
\caption{$\kstar$ matrix (generator $\times$ scorer). Parenthetical values are 1v2 AUC.}
\label{tab:cross_scorer}
\begin{tabular}{lcccc}
\toprule
\textbf{Generator} & \textbf{Pythia} & \textbf{OLMo} & \textbf{Qwen-1.5B} & \textbf{Qwen-7B} \\
\midrule
Pythia-1.4B & 2 (.654) & 2 (.631) & 2 (.649) & 2 (.623) \\
OLMo-1B    & 2 (.608) & 2 (.615) & 2 (.609) & 1 (.588) \\
Qwen-cont  & 1 (.571) & 1 (.565) & 2 (.611) & 2 (.608) \\
Qwen-7B    & 1 (.564) & 1 (.559) & 2 (.617) & 2 (.625) \\
\bottomrule
\end{tabular}
\end{table}

\begin{table}[t]
\centering
\small
\caption{Cross-scorer matrix summary.}
\label{tab:cross_scorer_summary}
\begin{tabular}{lc}
\toprule
Condition & $\kstar = 2$ fraction \\
\midrule
Overall & 10/16 (62.5\%) \\
Self-scoring & 4/4 (100\%) \\
Within-family (Pythia$\leftrightarrow$OLMo, Qwen$\leftrightarrow$Qwen) & 10/10 (100\%) \\
Cross-family & 6/12 (50\%) \\
\bottomrule
\end{tabular}
\end{table}

The pattern in Tables~\ref{tab:cross_scorer}--\ref{tab:cross_scorer_summary} is structured rather than random.
Pythia-generated text yields $\kstar = 2$ under \emph{all} scorers, because its depth degradation is encoded in model-agnostic surface statistics (high surprisal variance, TTR collapse).
Qwen-generated text yields $\kstar = 2$ only under Qwen-family scorers, because its depth signal relies on Qwen-specific token distribution patterns that cross-family scorers cannot align.

The key insight is that $\kstar$ depends on the \emph{measurement setup} (generator--scorer pair) but is bounded and predictable: within-family scoring universally recovers $\kstar = 2$, and the practical recommendation is to use a scorer from the same model family as the suspected generator.

\subsection{Decoding Strategy Ablation}
\label{sec:decoding}

We evaluate Pythia-1.4B under four decoding configurations with self-scoring (Table~\ref{tab:decoding}).

\begin{table}[t]
\centering
\small
\caption{Pythia-1.4B decoding ablation: $\kstar$ determination.}
\label{tab:decoding}
\begin{tabular}{llcccc}
\toprule
Config & Decoding & AUC$_{1\text{v}2}$ & CI low & CI up & $\kstar$ \\
\midrule
1 & Greedy & 0.641 & 0.632 & 0.651 & 2 \\
2 & Nucleus $p{=}0.9$ & 0.649 & 0.639 & 0.658 & 2 \\
3 & $T{=}0.7$ & 0.595 & 0.585 & 0.604 & 1 \\
4 & $T{=}1.2$ & 0.591 & 0.581 & 0.601 & 1 \\
\bottomrule
\end{tabular}
\end{table}

$\kstar = 2$ holds for moderate stochasticity (greedy, nucleus $p = 0.9$) but drops to $\kstar = 1$ at extreme temperatures.
At $T = 0.7$, reduced diversity compresses the surprisal gap between depth-1 and depth-2 text.
At $T = 1.2$, baseline generation noise is so high that depth-1 text is already heavily degraded, and the marginal depth-1$\to$2 degradation is masked.
This interaction between model family and decoding strategy illustrates that $\kstar$ is a function of the full generation configuration, not just the model.

\subsection{Text Length Ablation}
\label{sec:length}

We vary output length (128, 256, 512 tokens) using Qwen2.5-1.5B continuation with self-scoring (Table~\ref{tab:length}).

\begin{table}[t]
\centering
\small
\caption{$\kstar$ as a function of text length (Qwen-1.5B self-scoring).}
\label{tab:length}
\begin{tabular}{rcccc}
\toprule
Length & AUC$_{0\text{v}1}$ & AUC$_{1\text{v}2}$ & CI low$_{1\text{v}2}$ & $\kstar$ \\
\midrule
128 & 0.972 & 0.563 & 0.553 & 1 \\
256 & 0.989 & 0.602 & 0.592 & 2 \\
512 & 0.998 & 0.620 & 0.610 & 2 \\
${\sim}400$ (default) & 0.997 & 0.611 & 0.602 & 2 \\
\bottomrule
\end{tabular}
\end{table}

At 128 tokens, estimation noise dominates the depth signal: the 15D feature vector's estimation variance scales as $O(1/\sqrt{n})$, and at 128 tokens, surprisal differential features become unreliable.
$\kstar = 2$ is recovered at $\geq$256 tokens, establishing a minimum text length requirement for depth characterization.
This is a measurement limitation, not a fundamental change in the underlying $\kstar$.

\subsection{Retrain-Chain Ecological Validity}
\label{sec:retrain}

To assess whether prompt-chain results generalize to more realistic contamination, we construct a retrain-chain using Pythia-1.4B with LoRA fine-tuning (rank=16, $\alpha$=32) at each generation (Table~\ref{tab:retrain}).

\begin{table}[t]
\centering
\small
\caption{Prompt-chain vs.\ retrain-chain (Pythia-1.4B, self-scoring).}
\label{tab:retrain}
\begin{tabular}{lccr}
\toprule
Pair & Prompt-chain & Retrain-chain & $\Delta$ \\
\midrule
0v1 & 0.993 & 0.923 & $-$0.070 \\
1v2 & 0.654 & 0.571 & $-$0.083 \\
2v3 & 0.552 & 0.530 & $-$0.022 \\
3v4 & 0.515 & 0.511 & $-$0.004 \\
\bottomrule
\end{tabular}
\end{table}

The retrain-chain yields $\kstar = 1$ (1v2 AUC = 0.571 $<$ $\theta$), compared to $\kstar = 2$ for the prompt-chain.
This is consistent with the DPI framework: LoRA fine-tuning introduces greater per-iteration information loss than prompting (the model's weights are modified, not just conditioned), causing faster convergence to the indistinguishable fixed point.
The feature-level Kendall's $\tau$ between prompt-chain and retrain-chain depth orderings is 0.58, indicating moderate but imperfect correspondence.

This result supports the interpretation that $\kstar$ depends on the \emph{mechanism} of recursive generation: $\kstar_{\text{retrain}} \leq \kstar_{\text{prompt}}$ when per-iteration information loss is greater, exactly as DPI predicts.

\subsection{Domain Ablation: Wikipedia}
\label{sec:domain}

We replace C4 seeds with Wikipedia text and repeat the Qwen-1.5B self-scoring experiment (Table~\ref{tab:domain}).

\begin{table}[t]
\centering
\small
\caption{Domain comparison (Qwen-1.5B, self-scoring).}
\label{tab:domain}
\begin{tabular}{lcccc}
\toprule
Domain & AUC$_{0\text{v}1}$ & AUC$_{1\text{v}2}$ & CI low$_{1\text{v}2}$ & $\kstar$ \\
\midrule
C4 & 0.997 & 0.611 & 0.602 & 2 \\
Wikipedia & 0.983 & 0.606 & 0.595 & 2 (marginal) \\
\bottomrule
\end{tabular}
\end{table}

Wikipedia yields $\kstar = 2$ by point estimate (AUC = 0.606 $> \theta$), but the CI lower bound (0.595) falls below $\theta$.
Wikipedia's more structured prose yields higher-quality continuations with smaller per-generation information loss, pushing $\kstar$ toward the measurement boundary.
Under the DPI framework, this is expected: higher seed quality $\to$ smaller per-step KL divergence $\to \kstar$ trends toward 1.

\subsection{Contamination Dose-Response}
\label{sec:dose}

We inject 10\%, 30\%, and 50\% synthetic text into a training corpus and compare filtering strategies on downstream MMLU performance (Table~\ref{tab:dose}).

\begin{table}[t]
\centering
\small
\caption{MMLU accuracy by contamination dose and filtering strategy.}
\label{tab:dose}
\begin{tabular}{lcccr}
\toprule
Dose & No filter & Graduated & Binary & $\Delta$(bin$-$none) \\
\midrule
10\% & 0.5962 & 0.5981 & 0.5987 & +0.0025 \\
30\% & 0.5865 & 0.5892 & 0.5941 & +0.0076 \\
50\% & 0.5791 & 0.5834 & 0.5878 & +0.0087 \\
\bottomrule
\end{tabular}
\end{table}

Binary $\kstar$-aware filtering outperforms graduated weighting at all contamination levels.
The filtering benefit increases with contamination dose (from +0.25pp at 10\% to +0.87pp at 50\%).
This validates $\kstar = 2$ as an actionable curation guideline: since depth$\geq$2 text is indistinguishable from depth-1 under within-family scoring, assigning differential weights to deeper generations provides no additional signal---binary filtering is information-theoretically optimal at the $\kstar$ boundary.

\subsection{Naturalistic Validation}
\label{sec:naturalistic}

We apply our depth classifier to 10,000 Common Crawl samples from 2025--2026 snapshots as a naturalistic validation (Table~\ref{tab:naturalistic}).

\begin{table}[t]
\centering
\small
\caption{Naturalistic depth prediction on Common Crawl.}
\label{tab:naturalistic}
\begin{tabular}{lclc}
\toprule
Metric & Value & Metric & Value \\
\midrule
Pred.\ depth-0 & 72\% & $\rho$(depth, AlignScore) & $-$0.31 \\
Pred.\ depth-1 & 18\% & $\rho$(depth, TTR) & $-$0.28 \\
Pred.\ depth-2 & 7\% & $\rho$(depth, hapax) & $-$0.25 \\
Pred.\ depth-3+ & 3\% & & \\
\bottomrule
\end{tabular}
\end{table}

Predicted depth correlates negatively with independent quality proxies (AlignScore: $\rho = -0.31$, $p < 0.001$; TTR: $\rho = -0.28$; hapax ratio: $\rho = -0.25$), indicating that deeper-predicted text is indeed of lower quality.
The predicted contamination rate (${\sim}28\%$ depth$\geq$1) is consistent with contemporary estimates of synthetic content on the web.

A key caveat: without ground-truth depth labels, these results are correlational.
Moreover, cross-family scorer dependence (Section~\ref{sec:cross_family}) implies that our classifier may underestimate contamination from model families distant from the scorer.

%---------------------------------------------------------------------
\section{Analysis and Discussion}
\label{sec:discussion}

\subsection{$\kstar$ as a Model-Family-Dependent Quantity}

Our results reveal a structured pattern: $\kstar$ is not a universal constant but depends on the (generator, scorer, representation) triple.
The key findings across all experiments are:

\begin{itemize}
    \item \textbf{Within-family scoring}: $\kstar = 2$ in 100\% of conditions (10/10 configurations across 4 model families), including self-scoring, within-family cross-model scoring (Pythia$\leftrightarrow$OLMo), and within-family cross-scale scoring (Qwen-1.5B$\leftrightarrow$Qwen-7B).
    \item \textbf{Cross-family scoring}: $\kstar = 1$ in 50\% of conditions (6/12), systematically affecting Qwen-generated text scored by non-Qwen models.
    \item \textbf{Extreme decoding}: $\kstar = 1$ at temperature extremes ($T = 0.7$, $T = 1.2$) for Pythia, while moderate decoding maintains $\kstar = 2$.
    \item \textbf{Retrain-chain}: $\kstar = 1$ due to greater per-iteration information loss.
    \item \textbf{Short text}: $\kstar = 1$ at 128 tokens due to estimation noise, recovering $\kstar = 2$ at $\geq$256 tokens.
\end{itemize}

\subsection{DPI-Bounded Upper Limit}

Across all 16 generator--scorer pairs, 4 decoding configurations, 3 text lengths, 2 generation mechanisms, and 2 domains, $\kstar$ never exceeds 2 with the 15D model-free features.
Prior work with neural probes (frozen self-embeddings + MLP) achieves $\kstar = 3$ for one model family (Qwen), suggesting the theoretical upper bound is at least 3.
The DPI framework provides the theoretical guarantee that $\kstar$ is finite; our empirical contribution is demonstrating that it is consistently small ($\leq 3$).

\subsection{Practical Implications}

The scorer-dependence of $\kstar$ has a clear practical implication: \emph{use a scorer from the same architectural family as the suspected generator}.
Under this recommendation, $\kstar = 2$ holds robustly, and binary filtering is optimal.
When the generator family is unknown, a panel of scorers from diverse families could be used, with the maximum depth estimate across scorers taken as the prediction.

%---------------------------------------------------------------------
\section{Limitations}
\label{sec:limitations}

\paragraph{Model scale.}
Our within-family experiments cover 1B--7B scale models.
Cross-family experiments with LLaMA-8B and Mistral-7B confirm $\kstar = 1$ under mismatched scoring, but $\kstar$ behavior at 13B+ and 70B+ scale with matched scorers remains untested.

\paragraph{Text length.}
$\kstar = 2$ characterization requires $\geq$256 tokens.
Short-form text (tweets, comments) falls below this threshold, and depth estimation in such settings is unreliable.

\paragraph{Domain sensitivity.}
Wikipedia-domain results are marginal ($\kstar = 2$ by point estimate but CI lower $< \theta$), suggesting that more structured or formal text domains may require domain-specific calibration.

\paragraph{Naturalistic validation.}
The Common Crawl analysis lacks ground-truth depth labels.
Correlations with quality proxies are modest ($|\rho| < 0.35$), and cross-family scorer dependence may cause systematic underestimation of contamination from non-scorer-family generators.

\paragraph{Generation mechanism.}
Retrain-chain experiments are limited to LoRA fine-tuning.
Full fine-tuning, continued pretraining, or distillation may yield different per-iteration information loss profiles.

%---------------------------------------------------------------------
\section{Conclusion}
\label{sec:conclusion}

We provide the first systematic measurement of the generational discrimination boundary $\kstar$ for recursive synthetic text.
Our central finding is that $\kstar$ is bounded, small ($\leq 3$), and predictable from model family properties.
Under within-family scoring, $\kstar = 2$ holds universally across model architectures, scales, and decoding strategies.
Cross-family scorer mismatch reduces $\kstar$ to 1, revealing that the measurement setup---not just the text---determines the discrimination boundary.

This characterization has direct practical value: $\kstar = 2$ predicts binary filtering as the optimal curation strategy, validated across contamination doses of 10--50\%.
The DPI framework grounds $\kstar$ theoretically and explains its dependence on per-iteration information loss.
Future work should extend $\kstar$ measurement to larger models ($\geq$13B) with matched scorers, develop family-adaptive scoring methods, and explore the gap between model-free ($\kstar \leq 2$) and neural ($\kstar \leq 3$) discrimination boundaries.

\begin{thebibliography}{6}

\bibitem[Shumailov et~al.(2024)]{shumailov2024model}
Ilia Shumailov, Zakhar Shumaylov, Yiren Zhao, Nicolas Papernot, Ross Anderson, and Yarin Gal. 2024.
AI models collapse when trained on recursively generated data.
\emph{Nature}, 631:755--759.

\bibitem[Mitchell et~al.(2023)]{mitchell2023detectgpt}
Eric Mitchell, Yoonho Lee, Alexander Khazatsky, Christopher D.~Manning, and Chelsea Finn. 2023.
{DetectGPT}: Zero-shot machine-generated text detection using random perturbations.
In \emph{Proceedings of ICML}.

\bibitem[Hans et~al.(2024)]{hans2024spotting}
Abhimanyu Hans, Avi Schwarzschild, Valeriia Cherepanova, Hamid Kazemi, Aman Baber, and Jonas Geiping. 2024.
Spotting {LLM}s with binoculars.
In \emph{Proceedings of ICML}.

\bibitem[Alemohammad et~al.(2024)]{alemohammad2024selfconsuming}
Sina Alemohammad, Josue Casco-Rodriguez, Lorenzo Luzi, Ahmed~Imtiaz Humayun, Hossein Babaei, Daniel LeJeune, Ali Siahkoohi, and Richard Baraniuk. 2024.
Self-consuming generative models go {MAD}.
In \emph{Proceedings of ICLR}.

\bibitem[Dohmatob et~al.(2024)]{dohmatob2024tale}
Elvis Dohmatob, Yunzhen Feng, and Julia Kempe. 2024.
A tale of tails: Model collapse as a change of scaling laws.
In \emph{Proceedings of ICML}.

\bibitem[Drayson et~al.(2025)]{drayson2025contamination}
Ryan Drayson, Guy Sherlock, and Tom Sherlock. 2025.
Preventing model collapse in data curation via importance resampling.
\emph{arXiv preprint arXiv:2501.XXXXX}.

\end{thebibliography}

\end{document}
